\chapter{Pruebas y validación}
\label{chap:pruebas}

\ph{Este capítulo describe cómo se ha verificado que el sistema cumple los
requisitos. Las secciones siguen el orden clásico de la \emph{pirámide de
pruebas}: de la base (muchas pruebas, rápidas y baratas) a la cumbre (pocas,
lentas y caras). En la base están las pruebas unitarias, en el medio las
de integración y en la cima las de extremo a extremo / aceptación.
Delimita el alcance desde el principio: este capítulo responde a la pregunta
de la \emph{verificación} («¿hemos construido bien el sistema?»); la
\emph{evaluación} de los resultados frente a los objetivos del trabajo y su
discusión se realiza en el Capítulo~\ref{chap:resultados}.}

\section{Estrategia de pruebas}
\label{sec:estrategia_pruebas}

\ph{Resumen del enfoque adoptado: niveles de pruebas usados, framework
elegido y entorno en el que se ejecutan. Adapta los niveles al tipo de TFG
(software, hardware, modelo de IA, estudio empírico\dots).}

\ph{Niveles de la pirámide aplicados en este TFG (incluye los que apliquen):}

\ph{
\begin{itemize}
  \item \textbf{\ph{Pruebas unitarias}}: \ph{base de la pirámide. Verifican
        unidades aisladas (funciones, clases, módulos). Son las más numerosas
        y rápidas; deben ser deterministas y no depender de servicios externos.}
  \item \textbf{\ph{Pruebas de integración}}: \ph{nivel intermedio. Comprueban
        que varios componentes funcionan juntos correctamente, incluyendo
        contratos con servicios externos (BD, APIs internas\dots).}
  \item \textbf{\ph{Pruebas de extremo a extremo / sistema}}: \ph{cumbre de la
        pirámide. Validan flujos completos desde la perspectiva del usuario o
        cliente. Son las menos numerosas porque son las más lentas y frágiles.}
  \item \textbf{\ph{Pruebas de aceptación (opcional)}} : \ph{cotejan el sistema
        frente a los requisitos del Capítulo~\ref{chap:requisitos}.}
\end{itemize}
}

\ph{Comenta brevemente el entorno de pruebas: si las dependencias externas
son reales, simuladas o efímeras (en memoria, contenedores desechables,
mocks\dots) y cómo se garantiza el aislamiento entre tests.}

La estrategia de pruebas aplicada tiene como objetivo verificar el correcto
funcionamiento del generador integrado en UVLHub, validando tanto los
componentes individuales como el flujo completo de configuración y generación.

Las pruebas desarrolladas se organizan en los siguientes niveles:

\begin{itemize}

    \item \textbf{Pruebas unitarias}: enfocadas en validar componentes aislados
    del sistema, incluyendo las reglas de validación, la construcción de
    configuraciones y la lógica interna del generador.

    \item \textbf{Pruebas de integración}: comprueban la interacción entre los
    distintos componentes del módulo \texttt{generator}, incluyendo la gestión
    del estado del \textit{}, la persistencia de parámetros y la generación de
    salidas.

    \item \textbf{Pruebas de sistema}: verifican el funcionamiento completo del
    flujo desde la perspectiva del usuario, incluyendo la interacción con el
    asistente de configuración mediante pruebas automatizadas.

    \item \textbf{Pruebas manuales}: se han realizado comprobaciones manuales del
    flujo completo de generación, verificando el comportamiento de la aplicación
    durante el desarrollo.

\end{itemize}

La ejecución de las pruebas automatizadas se realiza mediante la suite de tests
del proyecto, utilizando el entorno definido mediante Docker para garantizar la
disponibilidad de las dependencias necesarias y la reproducibilidad de los
resultados.

\section{Pruebas unitarias}
\label{sec:pruebas_unitarias}

\ph{Describe cómo se prueban los componentes centrales del sistema en
aislamiento (la lógica de negocio, el algoritmo principal o el módulo más
crítico del TFG). Indica qué grupos de casos cubre, no la lista exhaustiva.}

\ph{Grupos de casos típicos:}

\ph{
\begin{itemize}
  \item \textbf{\ph{Caso nominal}}: \ph{entrada válida que debe producir el
        resultado esperado.}
  \item \textbf{\ph{Entradas inválidas}}: \ph{datos mal formados, fuera de
        rango o con campos ausentes.}
  \item \textbf{\ph{Casos límite}}: \ph{valores en la frontera del dominio
        admisible.}
  \item \textbf{\ph{Reglas de negocio}}: \ph{combinaciones que activan
        restricciones específicas del dominio.}
\end{itemize}
}

\ph{Si la lógica del módulo no se reduce a un booleano, comenta también qué
información complementaria se valida (mensajes, códigos de error, listas de
campos afectados\dots).}

Las pruebas unitarias se han centrado en validar los componentes internos del
generador y la lógica asociada al asistente de configuración. Los principales
grupos de casos cubiertos son:

\begin{itemize}

    \item \textbf{Generador de modelos de características:}
    se valida la operación
    \texttt{GenerateFeatureModel.execute}, comprobando que los parámetros del
    objeto \texttt{FmgeneratorModel} se aplican correctamente durante la
    generación. Se verifican aspectos como el número de modelos generados, la
    reproducibilidad mediante semillas, los límites de características y
    profundidad del árbol, así como la generación de relaciones jerárquicas,
    atributos y restricciones según los niveles UVL habilitados.

    \item \textbf{Generación de restricciones y atributos:}
    se comprueba que la activación de los diferentes niveles del generador
    produce únicamente los elementos permitidos. Se validan restricciones
    booleanas, aritméticas y de tipo string, además de operadores agregados
    (\texttt{sum}, \texttt{avg}), funciones sobre cadenas y distribución de
    tipos de atributos. También se verifica que los atributos configurados para
    no participar en restricciones no aparecen en las expresiones generadas.

    \item \textbf{Validación del modelo de parámetros:}
    se prueban las comprobaciones realizadas por
    \texttt{FmgeneratorModel}, incluyendo la validación de distribuciones de
    probabilidad, rangos de generación y coherencia entre niveles activados.
    Además, se comprueba la correcta conversión desde el diccionario plano de
    parámetros utilizado por UVLHub hacia la estructura interna del generador.

    \item \textbf{Validadores del asistente de configuración:}
    se prueban los métodos de validación de cada paso del \textit{wizard}
    (\texttt{validate\_step1\_form} a
    \texttt{validate\_step5\_form}), comprobando casos como límites inválidos
    de modelos o características, distribuciones incorrectas, configuraciones
    incompatibles entre niveles y restricciones imposibles de generar.

    \item \textbf{Gestión del estado y construcción de parámetros:}
    se valida la persistencia temporal de la configuración almacenada en la
    sesión de Flask mediante las funciones de gestión del estado del \textit{wizard}.
    También se comprueba la normalización de distribuciones y la correcta
    propagación de los valores entre pasos mediante los constructores de
    parámetros.

\end{itemize}

\section{Pruebas de integración}
\label{sec:pruebas_integracion}

\ph{Describe las pruebas que ejercitan varios componentes a la vez.
Conviene separarlas según los componentes que combinan.}

\ph{
\begin{itemize}
  \item \ph{Pruebas sobre el contrato con la base de datos / capa de
        persistencia (ORM, esquemas, transacciones).}
  \item \ph{Pruebas sobre la capa de servicios o API
        (códigos de respuesta, contratos de entrada/salida, autenticación,
        control de acceso entre usuarios).}
  \item \ph{Pruebas sobre integraciones con servicios externos
        (proveedores de identidad, almacenamiento, APIs de terceros\dots).}
\end{itemize}
}

\ph{Documenta también las pruebas de casos negativos
(valores fuera de rango, restricciones violadas, datos incompletos\dots).}

Las pruebas de integración se han centrado en verificar la comunicación entre
los diferentes módulos que intervienen en el proceso de generación, desde la
configuración realizada mediante el \textit{wizard} de UVLHub hasta la obtención del
modelo UVL generado.

Los principales casos cubiertos son:

\begin{itemize}

    \item \textbf{Integración del \textit{wizard} con la capa de servicios:}
    se prueba el recorrido completo del asistente mediante las rutas
    \texttt{/generator/random/step1} a
    \texttt{/generator/random/step6}, verificando la interacción entre los
    controladores Flask y los métodos de \texttt{GeneratorWizardService}
    encargados de procesar cada etapa.

    \item \textbf{Integración de la configuración persistida con el generador:}
    se comprueba que la configuración almacenada durante el flujo del \textit{wizard} se
    recupera correctamente mediante el endpoint
    \texttt{/generator/random/params-json} y puede ser transformada en una
    instancia válida de \texttt{FmgeneratorModel}.

    \item \textbf{Integración del modelo de parámetros con la generación UVL:}
    se valida que una configuración completa procedente del \textit{wizard} es utilizada
    correctamente por \texttt{GenerateFeatureModel}, generando artefactos UVL
    que reflejan las opciones seleccionadas, como niveles UVL, cardinalidades,
    tipos de atributos, restricciones y opciones de nomenclatura.

    \item \textbf{Integración del flujo de navegación y estado del asistente:}
    se verifican los escenarios de navegación hacia atrás entre pasos,
    comprobando que los valores almacenados en la sesión Flask permanecen
    consistentes y que una configuración parcialmente modificada no provoca
    estados inválidos durante la generación.

\end{itemize}

\section{Pruebas de extremo a extremo}
\label{sec:pruebas_e2e}

\ph{Describe la validación del sistema desde la perspectiva del usuario
final. No tiene por qué ejecutarse en un navegador o entorno gráfico real:
muchas veces basta con automatizar el flujo a través de la interfaz pública
del sistema (HTTP, CLI, etc.). Estas pruebas son las menos numerosas pero
las que dan más confianza en que el sistema funciona como un todo.}

\ph{Plantilla de pasos del flujo principal validado:}

\ph{\begin{enumerate}
  \item \ph{Paso 1.}
  \item \ph{Paso 2.}
  \item \ph{Paso 3.}
  \item \ph{Paso 4.}
  \item \ph{Paso 5.}
\end{enumerate}
}

\ph{Casos complementarios habituales:}

\ph{\begin{itemize}
  \item \ph{Estado inicial vacío.}
  \item \ph{Repetición del flujo con varias instancias y comprobación del orden.}
  \item \ph{Acceso no autorizado o sesión inválida.}
  \item \ph{Recuperación tras un error controlado.}
\end{itemize}
}

Las pruebas de extremo a extremo se han implementado mediante Selenium
WebDriver, validando el comportamiento del sistema desde la interfaz web hasta
la ejecución del proceso de generación en el navegador.

El flujo principal validado es el siguiente:

\begin{enumerate}

    \item \textbf{Carga del generador:}
    se verifica la correcta inicialización de la interfaz mediante el acceso a
    la ruta \texttt{/generator/random/step1} y la disponibilidad de los
    elementos necesarios para comenzar la configuración.

    \item \textbf{Ejecución del proceso de generación:}
    se automatiza la interacción con el \textit{wizard} completo, comprobando la
    transición entre las diferentes vistas y la correcta respuesta de la
    interfaz ante las acciones del usuario.

    \item \textbf{Inicialización del entorno Pyodide:}
    se valida la carga del runtime JavaScript necesario para ejecutar el
    generador en el navegador antes de iniciar la generación.

    \item \textbf{Generación y visualización del modelo UVL:}
    se comprueba la ejecución de la generación desde la interfaz, verificando
    el estado final del proceso, la aparición del resultado generado y la
    visualización del contenido UVL obtenido.

\end{enumerate}

Además, se han probado escenarios adicionales relacionados con la interacción
del usuario:

\begin{itemize}

    \item Navegación entre pasos mediante los controles de la interfaz,
    verificando que el comportamiento visual coincide con el estado esperado.

    \item Activación y desactivación de opciones dependientes dentro del
    formulario, como niveles aritméticos, cardinalidades y opciones avanzadas.

    \item Ejecución del generador con configuraciones completas desde la
    interfaz, comprobando la disponibilidad de las acciones finales sobre el
    modelo generado.

\end{itemize}

\section{Cobertura y métricas de calidad}
\label{sec:cobertura}

La calidad de la solución se ha evaluado mediante pruebas automatizadas y
cobertura de líneas, utilizando \texttt{pytest} y \texttt{coverage.py}. Se ha
establecido un umbral mínimo del \textbf{80\,\%} sobre el código de producción.

Este umbral busca mantener una cobertura honesta. Las líneas no cubiertas
corresponden principalmente a comprobaciones defensivas, configuraciones
inválidas bloqueadas por el \textit{wizard} y rutas de error cuya ejecución
forzada no aportaría evidencia funcional relevante.

La medición excluye las dependencias externas, las cachés, los artefactos
generados, los ficheros auxiliares de carga y los propios tests. Aunque los
informes incluyen estos ficheros para facilitar el análisis, no se utilizan
para calcular la cobertura del código de producción.

\begin{table}[H]
\centering
\small
\setlength{\tabcolsep}{5pt}
\renewcommand{\arraystretch}{1.2}
\begin{tabularx}{\textwidth}{@{}
>{\raggedright\arraybackslash}X
>{\centering\arraybackslash}p{1.25cm}
>{\centering\arraybackslash}p{1.55cm}
>{\centering\arraybackslash}p{1.45cm}
@{}}
\toprule
\textbf{Fichero} &
\textbf{Líneas} &
\textbf{Sin cubrir} &
\textbf{Cobertura} \\
\midrule

\path{app/features/generator/__init__.py}
& 2 & 0 & 100\,\% \\

\path{app/features/generator/assets/js/}
\path{fmgen_wrapper.py}
& 49 & 0 & 100\,\% \\

\path{app/features/generator/routes.py}
& 214 & 6 & 97\,\% \\

\path{app/features/generator/wizard.py}
& 810 & 158 & 80\,\% \\

\path{fm_generator/src/.../models/models.py}
& 325 & 42 & 87\,\% \\

\path{fm_generator/src/.../operations/}
\path{generate_feature_model.py}
& 666 & 73 & 89\,\% \\

\midrule
\textbf{Código de producción de UVLHub}
& \textbf{1075} & \textbf{164} & \textbf{85\,\%} \\

\textbf{Código de producción de \texttt{fm\_generator}}
& \textbf{997} & \textbf{115} & \textbf{88\,\%} \\

\midrule
\textbf{Total de código de producción}
& \textbf{2072} & \textbf{279} & \textbf{87\,\%} \\

\bottomrule
\end{tabularx}
\caption{Cobertura de líneas del código de producción del generador.}
\label{tab:cobertura_ficheros}
\end{table}

El informe de UVLHub muestra una cobertura global del 91\,\%, pero dicho valor
incluye también los ficheros de prueba. Al excluirlos, el código de producción
del generador alcanza un 85\,\%. Del mismo modo, el informe del paquete
\texttt{fm\_generator} muestra un 91\,\% incluyendo sus tests, mientras que la
cobertura específica del código de producción es del 88\,\%.

Los ficheros con menor cobertura son \texttt{wizard.py}, con un 80\,\%, y
\texttt{models.py}, con un 87\,\%. A pesar de ello, los caminos principales de
validación, persistencia de parámetros, generación, serialización y lectura de
modelos se encuentran cubiertos. El motor
\texttt{GenerateFeatureModel}, responsable de construir los modelos, alcanza
un 89\,\%.

La suite completa está formada por 220 pruebas de UVLHub y 44 pruebas del
paquete \texttt{fm\_generator}, todas superadas en sus respectivas ejecuciones.
Las dos ejecuciones del paquete independiente no representan pruebas
adicionales, sino una repetición de las mismas pruebas, una de ellas con
medición de cobertura.

\begin{table}[H]
\centering
\small
\begin{tabular}{@{}P{7cm} r r@{}}
\toprule
\textbf{Conjunto de pruebas} & \textbf{N.º de pruebas} &
\textbf{Resultado} \\
\midrule
Pruebas de UVLHub &
220 & 220/220 pasan \\

Pruebas de \texttt{fm\_generator} &
44 & 44/44 pasan \\

\midrule
\textbf{Total de la suite} &
\textbf{264} & \textbf{264/264 pasan} \\
\bottomrule
\end{tabular}
\caption{Resumen de la suite automatizada de pruebas.}
\label{tab:resumen_suite}
\end{table}

La suite combina pruebas unitarias sobre validadores, constructores,
persistencia de parámetros y reglas del motor, con pruebas de integración del
flujo del \textit{wizard} y de la generación de modelos. Además, incluye
12 pruebas con Selenium para verificar los flujos principales desde la
interfaz web.

Destaca la prueba asociada al objetivo O1, que recorre 26 combinaciones válidas
de niveles principales y secundarios. Para cada combinación se generan tres
modelos, lo que supone 78 modelos en total. Cada modelo se serializa a UVL y
se vuelve a cargar mediante \texttt{UVLReader}. En las configuraciones
booleanas, además, el modelo se transforma a PySAT y se ejecuta la operación
de satisfacibilidad. Esta prueba comprueba la compatibilidad del resultado con
las herramientas que consumirán posteriormente los modelos generados.

La suite aporta evidencia para los objetivos de generación de modelos con
diferentes niveles de expresividad (O1), configuración del proceso de
generación (O2), integración con UVLHub y Flamapy (O3) y disponibilidad de una
batería automatizada de validación (O5). La relación detallada entre
requisitos y pruebas se recoge en la
Tabla~\ref{tab:trazabilidad_rf_tests}.

% --- Transición al siguiente capítulo ---
\ph{Cierra el capítulo con un párrafo puente de 2--4 líneas: qué queda
establecido aquí y qué pregunta abre el capítulo siguiente. Por ejemplo:
«Verificado que el sistema está bien construido y que cumple los requisitos
especificados, queda la pregunta de fondo: ¿resuelve el problema planteado y
en qué medida? El Capítulo~\ref{chap:resultados} evalúa los resultados
obtenidos frente a los criterios de éxito y los discute críticamente».}

La estrategia de pruebas desarrollada permite verificar el correcto
funcionamiento del sistema y comprobar el cumplimiento de los principales
requisitos definidos. Una vez validada la implementación, el siguiente capítulo
analiza los resultados obtenidos y evalúa en qué medida la solución desarrollada
alcanza los objetivos planteados inicialmente.